% ============================================================
%   CHAPITRE : Barycentre
%   Niveau   : Terminale S / Terminale C
%   Mazava Maths
% ============================================================
\chapter{Barycentre}
\begin{citationchapitre}
	« La géométrie est la science du raisonnement correct\\
	sur des figures incorrectement dessinées. »\\[0.4em]
	— Henri Poincaré
\end{citationchapitre}

\begin{objectifs}
	\begin{itemize}[label=\textbullet]
		\item Définir et construire le barycentre d'un système de points pondérés.
		\item Utiliser les propriétés d'associativité et de linéarité du barycentre.
		\item Déterminer et construire les lignes de niveau (ensembles de points).
		\item Appliquer le barycentre en géométrie analytique et aux centres remarquables.
	\end{itemize}
\end{objectifs}


% ============================================================
\section{Définition et existence du barycentre}
% ============================================================

\begin{definition}
	Soient $A_1, A_2, \ldots, A_n$ des points du plan et $\alpha_1, \alpha_2, \ldots, \alpha_n$
	des réels tels que $\alpha_1 + \alpha_2 + \cdots + \alpha_n \neq 0$.
	
	Il existe un unique point $G$ vérifiant :
	$$
	\alpha_1\,\vect{GA_1} + \alpha_2\,\vect{GA_2} + \cdots + \alpha_n\,\vect{GA_n} = \vect{0}.
	$$
	Ce point $G$ s'appelle le \textbf{barycentre} du système de points pondérés
	$$
	\{(A_1,\alpha_1)\,;\,(A_2,\alpha_2)\,;\,\ldots\,;\,(A_n,\alpha_n)\}.
	$$
	
	Les réels $\alpha_i$ s'appellent les \textbf{coefficients} (ou \textbf{poids}) du système.
\end{definition}

\begin{attention}
	Le barycentre n'existe que si la somme des coefficients est \textbf{non nulle}.
	Si $\alpha_1 + \alpha_2 + \cdots + \alpha_n = 0$, le système n'admet pas de barycentre.
\end{attention}

\begin{propriete}{Caractérisation vectorielle}{}
	$G$ est le barycentre de $\{(A_i, \alpha_i)\}$ si et seulement si pour tout point $M$
	du plan :
	$$
	\sum_{i=1}^{n} \alpha_i\,\vect{MA_i} = \left(\sum_{i=1}^{n} \alpha_i\right)\vect{MG}.
	$$
\end{propriete}

\begin{exemple}
	Soient $A$, $B$ deux points et $\alpha, \beta$ deux réels avec $\alpha + \beta \neq 0$.
	
	Le barycentre $G$ de $\{(A,\alpha)\,;\,(B,\beta)\}$ vérifie :
	$$
	\alpha\,\vect{GA} + \beta\,\vect{GB} = \vect{0}
	\iff \vect{AG} = \frac{\beta}{\alpha+\beta}\,\vect{AB}.
	$$
	$G$ divise $[AB]$ dans le rapport $\dfrac{\beta}{\alpha}$ à partir de $A$.
\end{exemple}

\begin{propriete}{Cas de deux points}{}
	Le barycentre $G$ de $\{(A,\alpha)\,;\,(B,\beta)\}$ avec $\alpha + \beta \neq 0$
	vérifie :
	$$
	\vect{AG} = \frac{\beta}{\alpha+\beta}\,\vect{AB}.
	$$
	\begin{itemize}[label=\textbullet]
		\item Si $\alpha = \beta$ : $G$ est le \textbf{milieu} de $[AB]$.
		\item Si $\alpha + \beta > 0$ et $\alpha, \beta > 0$ : $G$ est entre $A$ et $B$.
		\item Si les coefficients sont de signes opposés : $G$ est à l'extérieur de $[AB]$.
	\end{itemize}
\end{propriete}


% ============================================================
\section{Propriétés fondamentales}
% ============================================================

\begin{propriete}{Associativité du barycentre}{}
	Soient $G_1$ le barycentre de $\{(A_i, \alpha_i)_{i \in I}\}$ et
	$G_2$ le barycentre de $\{(A_j, \alpha_j)_{j \in J}\}$, avec $I \cap J = \emptyset$.
	
	Alors le barycentre de l'ensemble $\{(A_i, \alpha_i)_{i \in I \cup J}\}$
	est aussi le barycentre de
	$$
	\left\{\left(G_1,\ \sum_{i \in I}\alpha_i\right)\,;\,\left(G_2,\ \sum_{j \in J}\alpha_j\right)\right\}.
	$$
\end{propriete}

\remarque
Cette propriété permet de calculer un barycentre \textbf{en plusieurs étapes},
en regroupant d'abord certains points entre eux.

\begin{exemple}
	Calculer le barycentre $G$ de $\{(A,1)\,;\,(B,2)\,;\,(C,3)\}$.
	
	\textbf{Étape 1 :} On calcule $G_1 = $ barycentre de $\{(A,1)\,;\,(B,2)\}$
	avec poids total $3$.
	
	\textbf{Étape 2 :} Par associativité, $G = $ barycentre de $\{(G_1, 3)\,;\,(C,3)\}$,
	donc $G$ est le \textbf{milieu} de $[G_1C]$.
\end{exemple}

\begin{propriete}{Linéarité — relation de Chasles généralisée}{}
	Pour tout point $M$ du plan et tout système $\{(A_i, \alpha_i)\}$ de somme
	$\sigma = \sum \alpha_i \neq 0$ de barycentre $G$ :
	$$
	\sum_{i=1}^{n} \alpha_i\,\vect{MA_i} = \sigma\,\vect{MG}.
	$$
	En particulier, pour tout vecteur $\vec{u}$ :
	$$
	\sum_{i=1}^{n} \alpha_i\,\overrightarrow{A_i} = \sigma\,\overrightarrow{G},
	$$
	où les vecteurs positions sont pris depuis une même origine.
\end{propriete}

\begin{propriete}{Coordonnées du barycentre}{}
	Dans un repère, si $A_i$ a pour coordonnées $(x_i, y_i)$, alors le barycentre
	$G$ de $\{(A_i, \alpha_i)\}$ a pour coordonnées :
	$$
	x_G = \frac{\sum \alpha_i x_i}{\sum \alpha_i},
	\qquad
	y_G = \frac{\sum \alpha_i y_i}{\sum \alpha_i}.
	$$
\end{propriete}

\begin{exemple}
	Soient $A(1\,;\,2)$, $B(3\,;\,0)$, $C(-1\,;\,4)$ avec coefficients $2$, $1$, $3$.
	
	$$
	x_G = \frac{2 \times 1 + 1 \times 3 + 3 \times (-1)}{2+1+3}
	= \frac{2 + 3 - 3}{6} = \frac{2}{6} = \frac{1}{3}.
	$$
	$$
	y_G = \frac{2 \times 2 + 1 \times 0 + 3 \times 4}{6}
	= \frac{4 + 0 + 12}{6} = \frac{16}{6} = \frac{8}{3}.
	$$
	
	Donc $G\!\left(\dfrac{1}{3}\,;\,\dfrac{8}{3}\right)$.
\end{exemple}


% ============================================================
\section{Centres remarquables d'un triangle}
% ============================================================

\begin{propriete}{Centres remarquables comme barycentres}{}
	Soit $ABC$ un triangle.
	\begin{enumerate}
		\item Le \textbf{centre de gravité} $G$ est le barycentre de
		$\{(A,1)\,;\,(B,1)\,;\,(C,1)\}$. Il vérifie $\vect{GA} + \vect{GB} + \vect{GC} = \vect{0}$.

		\item L'\textbf{orthocentre} $H$ est le barycentre de
		$$
		\{(A,\tan\hat{A})\,;\,(B,\tan\hat{B})\,;\,(C,\tan\hat{C})\}.
		$$
		
		\item Le \textbf{centre du cercle circonscrit} $O$ est le barycentre de
		$$
		\{(A,\sin 2\hat{A})\,;\,(B,\sin 2\hat{B})\,;\,(C,\sin 2\hat{C})\}.
		$$
	\end{enumerate}
\end{propriete}

\begin{propriete}{Relation d'Euler}{}
	Dans tout triangle $ABC$, en notant $G$ le centre de gravité,
	$H$ l'orthocentre et $O$ le centre du cercle circonscrit :
	$$
	\vect{OH} = 3\,\vect{OG}.
	$$
	Les points $O$, $G$, $H$ sont \textbf{alignés} (droite d'Euler) et $OG = \dfrac{1}{3}OH$.
\end{propriete}


% ============================================================
\section{Lignes de niveau}
% ============================================================

\begin{definition}
	\textbf{Ligne de niveau}\par
	Soit $\{(A_i, \alpha_i)\}$ un système de points pondérés de barycentre $G$.
	Pour tout réel $k$, l'\textbf{ensemble de niveau $k$} est :
	$$
	\mathscr{E}_k = \left\{ M \in \mathscr{P} \;\middle|\;
	\sum_{i=1}^{n} \alpha_i\,MA_i^2 = k \right\}.
	$$
\end{definition}

\begin{theoreme}
	\textbf{Réduction de la somme pondérée des carrés}\par
	Soit $G$ le barycentre de $\{(A_i, \alpha_i)\}$ de somme $\sigma = \sum \alpha_i \neq 0$.
	Pour tout point $M$ :
	$$
	\sum_{i=1}^{n} \alpha_i\,MA_i^2
	= \sigma\,MG^2 + \sum_{i=1}^{n} \alpha_i\,GA_i^2.
	$$
\end{theoreme}

\remarque
On pose $\displaystyle K = \sum_{i=1}^{n} \alpha_i\,GA_i^2$ (constante indépendante de $M$).
La formule devient :
$$
\sum_{i=1}^{n} \alpha_i\,MA_i^2 = \sigma\,MG^2 + K.
$$

\begin{propriete}{Nature des lignes de niveau}{}
	Avec les notations ci-dessus :
	
	\begin{enumerate}
		\item \textbf{Si $\sigma > 0$} : l'ensemble $\mathscr{E}_k$ est
		\begin{itemize}[label=\textbullet]
			\item un \textbf{cercle} de centre $G$ et de rayon $r = \sqrt{\dfrac{k-K}{\sigma}}$
			si $k > K$,
			\item le \textbf{point} $G$ si $k = K$,
			\item l'\textbf{ensemble vide} si $k < K$.
		\end{itemize}
		
		\item \textbf{Si $\sigma < 0$} : l'ensemble $\mathscr{E}_k$ est
		\begin{itemize}[label=\textbullet]
			\item un \textbf{cercle} de centre $G$ et de rayon $r = \sqrt{\dfrac{k-K}{\sigma}}$
			si $k < K$,
			\item le \textbf{point} $G$ si $k = K$,
			\item l'\textbf{ensemble vide} si $k > K$.
		\end{itemize}
	\end{enumerate}
\end{propriete}

\begin{astucebac}
	Pour déterminer la nature d'un ensemble de points défini par
	$\alpha\,MA^2 + \beta\,MB^2 = k$ :
	\begin{enumerate}
		\item Calculer le barycentre $G$ de $\{(A,\alpha)\,;\,(B,\beta)\}$.
		\item Appliquer la formule de réduction : $\alpha\,MA^2 + \beta\,MB^2 = (\alpha+\beta)\,MG^2 + K$.
		\item Conclure selon le signe de $\sigma = \alpha + \beta$ et la valeur de $k - K$.
	\end{enumerate}
\end{astucebac}

\begin{exemple}
	Déterminer l'ensemble des points $M$ tels que $MA^2 + 2\,MB^2 = 9$,
	où $A(0\,;\,0)$ et $B(3\,;\,0)$.
	
	\textbf{Étape 1 :} $\sigma = 1 + 2 = 3$. Le barycentre $G$ de $\{(A,1)\,;\,(B,2)\}$ vérifie
	$\vect{AG} = \dfrac{2}{3}\,\vect{AB}$, donc $G(2\,;\,0)$.
	
	\textbf{Étape 2 :} $K = 1 \times GA^2 + 2 \times GB^2 = 4 + 2 \times 1 = 6$.
	
	\textbf{Étape 3 :} $MA^2 + 2\,MB^2 = 3\,MG^2 + 6 = 9 \iff MG^2 = 1 \iff MG = 1$.
	
	L'ensemble est le \textbf{cercle} de centre $G(2\,;\,0)$ et de rayon $1$.
\end{exemple}


% ============================================================
\section{Exercices résolus}
% ============================================================

\exerciceresolu
\textbf{Construire un barycentre par associativité}

Soient $A$, $B$, $C$, $D$ quatre points. Construire le barycentre $G$ de
$$
\{(A,3)\,;\,(B,-1)\,;\,(C,2)\,;\,(D,2)\}.
$$

\solution

$\sigma = 3 + (-1) + 2 + 2 = 6 \neq 0$, donc $G$ existe.

\textbf{Étape 1 :} $G_1 =$ barycentre de $\{(A,3)\,;\,(B,-1)\}$, poids total $2$.
$$
\vect{AG_1} = \frac{-1}{2}\,\vect{AB}, \quad \text{donc } G_1 \text{ est à l'extérieur de } [AB].
$$

\textbf{Étape 2 :} $G_2 =$ barycentre de $\{(C,2)\,;\,(D,2)\}$, poids total $4$.
$$
G_2 = \text{milieu de } [CD].
$$

\textbf{Étape 3 :} Par associativité, $G =$ barycentre de $\{(G_1, 2)\,;\,(G_2, 4)\}$ :
$$
\vect{G_1 G} = \frac{4}{6}\,\vect{G_1 G_2} = \frac{2}{3}\,\vect{G_1 G_2}.
$$


\exerciceresolu
\textbf{Démontrer que trois points sont alignés par barycentre}

Soient $A$, $B$, $C$ trois sommets d'un triangle et $I$, $J$, $K$ les milieux
respectifs de $[BC]$, $[CA]$, $[AB]$.

Montrer que les médianes $[AI]$, $[BJ]$, $[CK]$ sont concourantes en le centre
de gravité $G$ de $\{(A,1)\,;\,(B,1)\,;\,(C,1)\}$.

\solution

Soit $G$ le barycentre de $\{(A,1)\,;\,(B,1)\,;\,(C,1)\}$. On a :
$$
\vect{GA} + \vect{GB} + \vect{GC} = \vect{0}.
$$

$I$ est le milieu de $[BC]$, donc $\vect{GI} = \dfrac{\vect{GB}+\vect{GC}}{2} = -\dfrac{\vect{GA}}{2}$.

Ainsi $\vect{GA} = -2\,\vect{GI}$, ce qui signifie que $G$, $A$ et $I$ sont alignés
(avec $AG = 2\,GI$), donc $G \in [AI]$.

Par le même raisonnement, $G \in [BJ]$ et $G \in [CK]$.

Les trois médianes se coupent donc en $G$, le centre de gravité.


\exerciceresolu
\textbf{Déterminer un ensemble de points — ligne de niveau}

Dans un repère orthonormé, on donne $A(-2\,;\,1)$ et $B(4\,;\,3)$.

Déterminer et construire l'ensemble $\mathscr{E}$ des points $M$ tels que :
$$
2\,MA^2 - MB^2 = 10.
$$

\solution

\textbf{Étape 1 — Barycentre.}
$\sigma = 2 + (-1) = 1 \neq 0$.

Le barycentre $G$ de $\{(A,2)\,;\,(B,-1)\}$ vérifie :
$$
\vect{AG} = \frac{-1}{1}\,\vect{AB} = -\vect{AB}.
$$
Donc $x_G = x_A + (-1)(x_B - x_A) = -2 + (-1)(6) = -8$ et
$y_G = 1 + (-1)(2) = -1$.

Ainsi $G(-8\,;\,-1)$.

\textbf{Étape 2 — Constante $K$.}
$$
K = 2\,GA^2 + (-1)\,GB^2.
$$
$GA^2 = (-2-(-8))^2 + (1-(-1))^2 = 36 + 4 = 40$.

$GB^2 = (4-(-8))^2 + (3-(-1))^2 = 144 + 16 = 160$.

$K = 2 \times 40 - 160 = 80 - 160 = -80$.

\textbf{Étape 3 — Conclusion.}
$$
2\,MA^2 - MB^2 = 1 \times MG^2 + (-80) = 10
\iff MG^2 = 90
\iff MG = 3\sqrt{10}.
$$

$\mathscr{E}$ est le \textbf{cercle} de centre $G(-8\,;\,-1)$ et de rayon $3\sqrt{10}$.


\exerciceresolu
\textbf{Barycentre et vecteur nul}

Soient $A$, $B$, $C$ trois points non alignés. On cherche les réels $\alpha$, $\beta$, $\gamma$
tels que $\alpha + \beta + \gamma = 0$ et
$$
\alpha\,\vect{MA} + \beta\,\vect{MB} + \gamma\,\vect{MC} = \vect{u}
$$
pour tout point $M$, où $\vec{u}$ est un vecteur donné.

\solution

Lorsque $\alpha + \beta + \gamma = 0$, on sait que :
$$
\alpha\,\vect{MA} + \beta\,\vect{MB} + \gamma\,\vect{MC}
$$
est \textbf{indépendant du point} $M$ (c'est une combinaison libre de $\vect{AB}$ et $\vect{AC}$).

En développant depuis un point $O$ quelconque :
$$
\alpha\,\vect{OA} + \beta\,\vect{OB} + \gamma\,\vect{OC} = \vec{u}.
$$

Cela revient à décomposer $\vec{u}$ dans la base $(\vect{AB}, \vect{AC})$, ce qui admet
une solution unique puisque $A$, $B$, $C$ sont non alignés.

\begin{attention}
	Quand la somme des coefficients est nulle, la combinaison
	$\sum \alpha_i\,\vect{MA_i}$ ne dépend pas de $M$ : c'est un vecteur fixe,
	pas forcément $\vect{0}$. Il ne faut pas confondre avec la définition du barycentre.
\end{attention}


% ============================================================
\section{Exercices}
% ============================================================

\exercice{Barycentre de deux points \hfill \etoilepleine\etoilepleine\etoilevide}

Soient $A(1\,;\,3)$ et $B(5\,;\,-1)$.

\begin{enumerate}
	\item Déterminer le barycentre $G_1$ de $\{(A,3)\,;\,(B,1)\}$.
	\item Déterminer le barycentre $G_2$ de $\{(A,2)\,;\,(B,-2)\}$.
	Ce point existe-t-il ? Justifier.
	\item Déterminer le barycentre $G_3$ de $\{(A,1)\,;\,(B,4)\}$
	et vérifier que $G_3$ est bien entre $A$ et $B$.
\end{enumerate}


\exercice{Barycentre de trois points et centre de gravité \hfill \etoilepleine\etoilepleine\etoilevide}

Soient $A(0\,;\,0)$, $B(6\,;\,0)$, $C(0\,;\,4)$.

\begin{enumerate}
	\item Calculer les coordonnées du centre de gravité $G$ du triangle $ABC$.
	\item Vérifier que $\vect{GA} + \vect{GB} + \vect{GC} = \vect{0}$.
	\item Déterminer le barycentre $H$ de $\{(A,4)\,;\,(B,1)\,;\,(C,-2)\}$.
	\item Montrer que $G$, $H$ et $C$ sont-ils alignés ?
\end{enumerate}


\exercice{Ligne de niveau — cercle ou vide \hfill \etoilepleine\etoilepleine\etoilevide}

On donne $A(1\,;\,0)$ et $B(-1\,;\,0)$.

Déterminer la nature de l'ensemble des points $M$ vérifiant :
\begin{enumerate}
	\item $MA^2 + MB^2 = 6$
	\item $MA^2 + MB^2 = 1$
	\item $3\,MA^2 + MB^2 = 12$
	\item $MA^2 - 3\,MB^2 = 4$
\end{enumerate}


\exercice{Problème de synthèse \hfill \etoilepleine\etoilepleine\etoilepleine}

Dans un repère orthonormé, on donne $A(2\,;\,1)$, $B(-1\,;\,3)$, $C(4\,;\,-2)$.

\begin{enumerate}
	\item Déterminer le barycentre $G$ de $\{(A,1)\,;\,(B,2)\,;\,(C,3)\}$.
	
	\item Montrer que pour tout point $M$ :
	$$
	MA^2 + 2\,MB^2 + 3\,MC^2 = 6\,MG^2 + K,
	$$
	où $K$ est une constante que l'on calculera.
	
	\item En déduire la nature et les éléments caractéristiques de l'ensemble $\mathscr{E}$
	des points $M$ tels que $MA^2 + 2\,MB^2 + 3\,MC^2 = 30$.
	
	\item Déterminer le point $M_0$ de l'axe des abscisses appartenant à $\mathscr{E}$.
	
	\item Montrer que parmi tous les points du plan, c'est $G$ qui minimise
	$MA^2 + 2\,MB^2 + 3\,MC^2$. Quelle est cette valeur minimale ?
\end{enumerate}